\documentclass[12pt, a4paper]{article}

\usepackage{amsmath, amssymb, amsfonts}
\usepackage{graphicx}
\usepackage{geometry}
\geometry{left=2.5cm, right=2.5cm, top=3cm, bottom=3cm}
\usepackage{hyperref}
\usepackage{booktabs}
\usepackage{float}
\usepackage{caption}
\usepackage{subcaption}

\title{\textbf{Sparse Optimization Algorithms Based on Lasso and Their Applications in Compressed Sensing}\\ 
\large Final Project Report for Machine Learning Numerical Analysis}
\author{}
\date{\today}

\begin{document}

\maketitle

\begin{abstract}
This report delves into the application of the Lasso model in Compressed Sensing (CS) problems, aiming to achieve precise reconstruction of sparse signals through $\ell_1$ regularization. Starting from mathematical principles, the report details the theoretical derivations and code implementations of four classic numerical optimization algorithms: Iterative Shrinkage-Thresholding Algorithm (ISTA), Fast Iterative Shrinkage-Thresholding Algorithm (FISTA), Coordinate Descent (CD), and Alternating Direction Method of Multipliers (ADMM). These algorithms are comprehensively evaluated under two core experimental scenarios: the accurate recovery of one-dimensional artificial random sparse signals, and the compressed sensing reconstruction of a two-dimensional Shepp-Logan medical image phantom (simulating MRI/CT). In the 2D image experiment, the Discrete Cosine Transform (DCT) is introduced as the sparsifying basis, and a "Variable Density Random Fourier Sampling" strategy is specifically designed to enforce the preservation of low-frequency central energy. Experimental results demonstrate that this sampling strategy significantly improves the goodness of fit ($R^2 > 0.95$) of the reconstructed images. In the comparison of the four solvers, the FISTA algorithm strikes the best balance between computational efficiency and convergence accuracy. While ADMM requires the fewest iterations to converge, it incurs a substantial computational overhead for matrix inversion in a single step. The baseline ISTA and CD algorithms are disadvantaged in terms of either time consumption or the number of iterations.

\textbf{Keywords}: Lasso, Compressed Sensing, Soft Thresholding Operator, ISTA, FISTA, ADMM, Medical Image Reconstruction
\end{abstract}

\newpage

\section{Introduction}
Compressed Sensing (CS) theory breaks through the limitations of the Nyquist sampling theorem, indicating that if a signal is sparse in a certain transform domain, it can be precisely reconstructed using far less data than conventionally required. This theory holds extremely high application value in fields such as medical imaging (e.g., rapid MRI scanning), radar detection, and image compression.

In compressed sensing, the reconstruction of a sparse signal is typically formulated as a convex optimization problem regularized by the $\ell_1$ norm, known as the Lasso (Least Absolute Shrinkage and Selection Operator) problem. Because the $\ell_1$ norm is non-differentiable at the origin, traditional gradient descent methods cannot be directly applied to solve it. Therefore, a series of efficient non-smooth optimization algorithms have emerged in the fields of machine learning and numerical analysis.

The core objective of this project is to independently implement four major numerical optimization algorithms: ISTA, FISTA, CD, and ADMM. Furthermore, these algorithms are applied to the reconstruction tasks of 1D signals and 2D medical images. Through a quantitative comparison of multiple dimensions (goodness of fit, MSE, runtime, convergence speed), we explore the performance boundaries of each algorithm in practical applications.

\section{Algorithm Mathematical Principles and Implementation}

In compressed sensing, the observation data $y \in \mathbb{R}^M$ and the unknown signal $x \in \mathbb{R}^N$ ($M < N$) satisfy the linear relationship $y = Ax + \epsilon$, where $A$ is the measurement matrix and $\epsilon$ is Gaussian noise. Given the sparsity of $x$, we reconstruct it by minimizing the following Lasso objective function:
\begin{equation}
    \min_{x} f(x) = \frac{1}{2} \|Ax - y\|_2^2 + \lambda \|x\|_1
\end{equation}

Due to the non-differentiability of the $\ell_1$ norm, we introduce the Soft Thresholding Operator as its Proximal Operator:
\begin{equation}
    S_{\kappa}(v) = \text{sgn}(v) \max(|v| - \kappa, 0)
\end{equation}

\subsection{Iterative Shrinkage-Thresholding Algorithm (ISTA)}
ISTA is a specific instance of the proximal gradient descent method applied to the Lasso problem. In each iteration step, gradient descent is first performed along the differentiable quadratic loss component, and then the soft thresholding operator is applied to handle the $\ell_1$ regularization term:
\begin{equation}
    x_{k+1} = S_{\alpha\lambda}\left( x_k - \alpha A^T(Ax_k - y) \right)
\end{equation}
where $\alpha \le 1/L$ is the step size, and $L = \|A\|_2^2$ is the Lipschitz constant of the smooth part of the objective function.

\subsection{Fast Iterative Shrinkage-Thresholding Algorithm (FISTA)}
To address the slow convergence speed of ISTA (convergence rate $O(1/k)$), FISTA introduces the Nesterov momentum acceleration mechanism, improving its convergence rate to $O(1/k^2)$. The iterative format is as follows:
\begin{align}
    x_{k+1} &= S_{\alpha\lambda}\left( z_k - \alpha A^T(Az_k - y) \right) \\
    t_{k+1} &= \frac{1 + \sqrt{1 + 4t_k^2}}{2} \\
    z_{k+1} &= x_{k+1} + \frac{t_k - 1}{t_{k+1}}(x_{k+1} - x_k)
\end{align}

\subsection{Coordinate Descent (CD)}
Coordinate descent minimizes the objective with respect to only one component of the variable $x$, denoted as $x_j$, at a time, while fixing the other components. For the Lasso problem, each coordinate subproblem has an analytical solution, which avoids the sensitivity to the step size in global gradient calculations. The formula is:
\begin{equation}
    x_j \leftarrow \frac{S_{\lambda}(A_j^T r + \|A_j\|_2^2 x_j)}{\|A_j\|_2^2}
\end{equation}
where $r = y - Ax$ is the current residual.

\subsection{Alternating Direction Method of Multipliers (ADMM)}
ADMM decomposes the original problem into independently solvable subproblems by introducing an auxiliary variable $z$. The constrained problem is written as: $\min \frac{1}{2}\|Ax - y\|_2^2 + \lambda\|z\|_1$, s.t. $x - z = 0$. Using the Augmented Lagrangian method, the update formulas are derived as:
\begin{align}
    x_{k+1} &= (A^TA + \rho I)^{-1}(A^Ty + \rho(z_k - u_k)) \\
    z_{k+1} &= S_{\lambda/\rho}(x_{k+1} + u_k) \\
    u_{k+1} &= u_k + x_{k+1} - z_{k+1}
\end{align}

\section{Experiment 1: 1D Sparse Signal Recovery}
\subsection{Experimental Setup}
We generated a ground truth signal $x_{true}$ with a dimension of $N=500$, a measurement dimension of $M=150$, and a sparsity of $S=10$. The measurement matrix $A$ adopts a Gaussian random matrix, and Gaussian noise with standard deviation $\sigma=0.05$ is added to the observations $y$. Four algorithms were used to reconstruct it.

\subsection{Results and Analysis}
Experimental results indicate that under identical initial conditions:
\begin{itemize}
    \item \textbf{Convergence Speed}: The objective function of ADMM drops extremely fast, converging in fewer than 60 steps. Utilizing momentum acceleration, FISTA converges in approximately 150 steps. ISTA is the slowest, requiring nearly 200 steps. While CD involves fewer iterations in terms of Epochs, each Epoch inherently contains $N$ internal loops.
    \item \textbf{Recovery Accuracy}: All four algorithms accurately pinpointed the support set of non-zero elements of the signal, and the reconstructed signals were highly consistent with the original signal at all peaks.
\end{itemize}

\section{Experiment 2: 2D Medical Image Compressed Reconstruction (MRI Simulation)}
\subsection{Variable Density Fourier Sampling and Sparsification Strategy}
In real Magnetic Resonance Imaging (MRI) applications, applying purely uniform random sampling directly to high-dimensional images results in extremely poor reconstruction quality. Therefore, we introduced the following critical improvement strategies:
\begin{enumerate}
    \item \textbf{Sparsifying Basis Selection}: The 2D Inverse Discrete Cosine Transform (IDCT) matrix was used as the transform basis $D$ to project the non-sparse Shepp-Logan phantom from the spatial domain to the frequency domain, rendering it highly sparse.
    \item \textbf{Variable Density Sampling}: The vast majority of MRI image energy is concentrated in the low-frequency region. When generating the sampling mask, we \textbf{enforced 100\% full sampling in the ultra-low frequency central region (radius $\le 6$)}, while for the peripheral high-frequency regions, the sampling probability decreased proportionally to the inverse cube of the distance. The overall sampling rate $M/N$ for the experiment was set to $60\%$.
    \item \textbf{Complex-to-Real Decoupling}: The matrix generated by Fourier sampling is complex-valued. To reuse our implemented real-domain Lasso solvers, the complex equations were decomposed and concatenated into real and imaginary parts: $A_{real} = [\Re(A)^T, \Im(A)^T]^T$.
\end{enumerate}

\subsection{Evaluation Metrics and Comparative Results}
We executed the reconstruction process using the four algorithms on an image of size $32 \times 32$, with a maximum of 2000 iterations and a tolerance of $10^{-5}$. The evaluation metrics were Mean Squared Error (MSE) and Goodness of Fit ($R^2$).

The variable density sampling scheme produced extremely clear reconstructed images, effectively suppressing artifacts. The $R^2$ scores obtained by all four algorithms successfully surpassed the $0.95$ threshold, proving the exceptional feasibility of the proposed scheme in compressed sensing image reconstruction.

\section{Comprehensive Evaluation of Algorithm Performance and Conclusion}
Combining the charts and data from the image and signal reconstructions, this report draws the following conclusions:
\begin{enumerate}
    \item \textbf{ADMM}: It holds an overwhelming advantage in terms of the "number of iterations," converging exceptionally fast. However, its update step involves computing the inverse of the matrix $A^TA + \rho I$. When the signal dimension is large, even with pre-computation and caching of the inverse matrix, the constant overhead of a single matrix multiplication remains enormous, resulting in suboptimal overall runtime.
    \item \textbf{CD (Coordinate Descent)}: Because it cannot perform global parallel matrix-vector multiplications, executing large-scale for-loops in an interpreted language like Python causes a severe performance bottleneck. Consequently, its computation time far exceeds the other algorithms.
    \item \textbf{ISTA}: It is the simplest to implement, requires no matrix inversion, and has extremely fast single steps. However, limited by the sublinear convergence rate of first-order gradients, it necessitates a large number of global iteration steps.
    \item \textbf{FISTA}: \textbf{Demonstrates the best overall performance.} Its single-step overhead is identical to that of ISTA (involving only matrix-vector multiplications), but leveraging Nesterov momentum drastically reduces the number of iterations. It showcases the most outstanding engineering practicality in trading off "computational time" against "reconstruction quality," making it the preferred solution for handling such large-scale $\ell_1$ optimization problems.
\end{enumerate}

In summary, this project comprehensively implemented various sparse optimization algorithms from theory to code, profoundly validating the immense potential of the Lasso model and variable density Fourier sampling strategies in the fields of modern compressed sensing and medical image reconstruction.

\end{document}
